% GENERATED FILE. Edit canonical lessons, then run npm run generate:books.
% canonical-source-hash: fnv1a64:7bb856ad72685757
% canonical-lessons: ZH-C157-kaihui, ZH-C157-tongzhi, ZH-C157-xiaoxi, ZH-C157-huifu, ZH-C157-reading, ZH-W157-reply, ZH-R157-review-final-157

\chapter{A Note About Tomorrow}
\label{ch:zh-a-note-about-tomorrow}
\hlchaptermodality{157}

\begin{chapteropening}
\textbf{By the end of this chapter:} \emph{I can understand a notice and a message, and write a connected reply.}

Complete ZH-R157-review-final-157: one dialogue across the A2 core, after reading a notice at work and a friend's message and writing a connected reply of four to six lines.
\end{chapteropening}

\section[\texorpdfstring{\zh{开会} (kāihuì)}{开会 (kāihuì)}]{\zh{开会} (kāihuì) — to have a meeting}

\label{lesson:ZH-C157-kaihui}



\begin{quote}
A new chapter, and it opens on a word you will hear at work.

\begin{itemize}
\raggedright
  \item \emph{Recall:} say \zh{一本书十元} — \emph{yì běn shū shí yuán} — \textbf{R1}, the chapter you have finished
  \item \emph{Recall:} say \emph{guǎnggào}, an advert — \textbf{R1}
  \item \emph{Recall:} say \zh{开门} — \emph{kāimén}, to open — \textbf{R2}
\end{itemize}
\end{quote}

\subsection*{You'll want to know: \zh{开会}}
\textbf{\zh{开会}} — \emph{kāihuì} — \textquotedblleft{}to have a meeting\textquotedblright{}.

\begin{quote}
\textbf{\zh{明天下午开会}} \emph{míngtiān xiàwǔ kāihuì} — There's a meeting tomorrow afternoon.
\end{quote}

\begin{quote}
\textbf{\zh{我们在开会}} \emph{wǒmen zài kāihuì} — We're in a meeting.
\end{quote}

\begin{cousinweb}
\textbf{\zh{开会}} (\emph{kāihuì}) is built from \zh{开} 'to open' and \zh{会} 'a gathering'. \zh{会} is the character of \zh{会} 'can' and \zh{会} 'will', inherited from Middle Chinese \emph{hwajH} (Guangyun, in Baxter's transcription), glossed 'to come together': here it keeps that sense.
\end{cousinweb}

\begin{grammarlens}[title={What you've built}]
\zh{开} starts many things you know: \zh{开门}, \zh{开车}, and now \zh{开会}.

\zh{在} before it says the meeting is going on: \zh{我们在开会}.
\end{grammarlens}

\subsection*{Your turn}
\begin{itemize}
\raggedright
  \item \emph{Say it:} \emph{kāihuì}
  \item \emph{Say it:} \emph{kāihuì}, once more
  \item \emph{Say it:} \zh{明天下午开会} — \emph{míngtiān xiàwǔ kāihuì} — There's a meeting tomorrow afternoon.
  \item \emph{Say it:} \zh{我们在开会} — \emph{wǒmen zài kāihuì} — We're in a meeting.
  \item \emph{Recall:} say \zh{开门} and \zh{开车}, then \zh{开会}
\end{itemize}

\begin{tcolorbox}[breakable,colback=teal!4,colframe=teal!35!black,title={Before you move on}]
What does \zh{开会} mean? (\textquotedblleft{}To have a meeting\textquotedblright{}.) Say there's a meeting tomorrow afternoon. (\textbf{\zh{明天下午开会}}.)
\end{tcolorbox}
\section[tōngzhī]{tōngzhī — to notify; a notice}

\label{lesson:ZH-C157-tongzhi}



\begin{quote}
Before the new one: last lesson's word, then a word that ends the same way.

\begin{itemize}
\raggedright
  \item \emph{Recall:} say \zh{开会} — \emph{kāihuì} — \textbf{R1}, one lesson back
  \item \emph{Recall:} say \emph{zhīdào}, to know — \textbf{R3}, durable
\end{itemize}
\end{quote}

\subsection*{You'll want to know: tōngzhī}
\emph{tōngzhī} — \textquotedblleft{}to notify; a notice\textquotedblright{}.

\begin{quote}
\emph{lǎoshī tōngzhī wǒmen míngtiān kāihuì} — The teacher let us know there's a meeting tomorrow.
\end{quote}

\begin{quote}
\emph{zhè shì yí ge tōngzhī} — This is a notice.
\end{quote}

Its characters come later; for now it is a word for the ear.

\begin{cousinweb}
\emph{tōngzhī} is built from tōng 'to get through' and zhī 'to know', the zhī of \emph{zhīdào}. Both are inherited from Middle Chinese, \emph{thuwng} and \emph{trje} (Guangyun, in Baxter's transcription), glossed 'to reach' and 'to perceive': making something known.
\end{cousinweb}

\begin{grammarlens}[title={What you've built}]
As a verb, \emph{tōngzhī} takes the people told, then the news: \emph{tōngzhī wǒmen míngtiān kāihuì}.

As a noun it is the notice itself, pinned up or sent round.
\end{grammarlens}

\subsection*{Your turn}
\begin{itemize}
\raggedright
  \item \emph{Say it:} \emph{tōngzhī}
  \item \emph{Say it:} \emph{tōngzhī}, once more
  \item \emph{Say it:} \emph{lǎoshī tōngzhī wǒmen míngtiān kāihuì} — The teacher let us know there's a meeting tomorrow.
  \item \emph{Recall:} say \emph{zhīdào}, then \emph{tōngzhī}
\end{itemize}

\begin{tcolorbox}[breakable,colback=teal!4,colframe=teal!35!black,title={Before you move on}]
What does \emph{tōngzhī} mean? (\textquotedblleft{}To notify; a notice\textquotedblright{}.) What comes after it as a verb? (\textbf{Who is told, then the news}.)
\end{tcolorbox}
\section[xiāoxi]{xiāoxi — news; a message}

\label{lesson:ZH-C157-xiaoxi}



\begin{quote}
Before the new one: last lesson's word, then a word that ends on the same light xi.

\begin{itemize}
\raggedright
  \item \emph{Recall:} say \emph{tōngzhī}, a notice — \textbf{R1}, one lesson back
  \item \emph{Recall:} say \emph{xiūxi}, to rest — \textbf{R3}, durable
\end{itemize}
\end{quote}

\subsection*{You'll want to know: xiāoxi}
\emph{xiāoxi} — \textquotedblleft{}news; a message\textquotedblright{}.

\begin{quote}
\emph{wǒ yǒu yí ge hǎo xiāoxi} — I've got some good news.
\end{quote}

Its characters come later; for now it is a word for the ear.

\begin{cousinweb}
\emph{xiāoxi} is built from xiāo and xi, the xi of \emph{xiūxi}. Both are inherited from Middle Chinese, \emph{sjew} and \emph{sik} (Guangyun, in Baxter's transcription), glossed 'to melt away; to wane' and 'to stop; to breed'. Neither gloss gives 'news', so how the pair came to mean it is not settled.
\end{cousinweb}

\begin{grammarlens}[title={What you've built}]
\emph{xiāoxi} is any news, and the message that brings it: a text from a friend is a \emph{xiāoxi}.

The xi at the end is light, as in \emph{xiūxi}.
\end{grammarlens}

\subsection*{Your turn}
\begin{itemize}
\raggedright
  \item \emph{Say it:} \emph{xiāoxi}
  \item \emph{Say it:} \emph{xiāoxi}, once more
  \item \emph{Say it:} \emph{wǒ yǒu yí ge hǎo xiāoxi} — I've got some good news.
  \item \emph{Recall:} say \emph{xiūxi}, then \emph{xiāoxi}
\end{itemize}

\begin{tcolorbox}[breakable,colback=teal!4,colframe=teal!35!black,title={Before you move on}]
What does \emph{xiāoxi} mean? (\textquotedblleft{}News; a message\textquotedblright{}.) Say you have some good news. (\emph{wǒ yǒu yí ge hǎo xiāoxi}.)
\end{tcolorbox}
\section[huífù]{huífù — to reply}

\label{lesson:ZH-C157-huifu}



\begin{quote}
Before the new one: last lesson's word, then two words this one leans on.

\begin{itemize}
\raggedright
  \item \emph{Recall:} say \emph{xiāoxi}, a message — \textbf{R1}, one lesson back
  \item \emph{Recall:} say \emph{huídá}, to answer — \textbf{R2}, a few chapters back
  \item \emph{Recall:} say \emph{fùxí}, to revise — \textbf{R2}
\end{itemize}
\end{quote}

\subsection*{You'll want to know: huífù}
\emph{huífù} — \textquotedblleft{}to reply (to a message)\textquotedblright{}.

\begin{quote}
\emph{qǐng huífù} — Please reply.
\end{quote}

Its characters come later; for now it is a word for the ear.

\begin{cousinweb}
\emph{huífù} is built from \emph{huí} 'to go back' and fù 'to return', the fù of \emph{fùxí}. Both are inherited from Middle Chinese, \emph{hwoj} and \emph{bjuwk} (Guangyun, in Baxter's transcription), glossed 'to turn round' and 'to return'. The regular change from \emph{bjuwk} gives fú; fù continues the Guangyun's other reading, \emph{bjuwH}, glossed 'again'.
\end{cousinweb}

\begin{grammarlens}[title={What you've built}]
\emph{huídá} answers a question; \emph{huífù} answers a message.

\zh{请} in front makes it the polite request at the foot of a notice: \emph{qǐng huífù}.
\end{grammarlens}

\subsection*{Your turn}
\begin{itemize}
\raggedright
  \item \emph{Say it:} \emph{huífù}
  \item \emph{Say it:} \emph{huífù}, once more
  \item \emph{Say it:} \emph{qǐng huífù} — Please reply.
  \item \emph{Recall:} say \emph{huídá}, then \emph{huífù}: a question, then a message
\end{itemize}

\begin{tcolorbox}[breakable,colback=teal!4,colframe=teal!35!black,title={Before you move on}]
What does \emph{huífù} mean? (\textquotedblleft{}To reply to a message\textquotedblright{}.) And \emph{huídá}? (\textbf{To answer a question}.)
\end{tcolorbox}
\section[(reading)]{(reading) — a notice and a message}

\label{lesson:ZH-C157-reading}



\begin{quote}
A notice pinned up at work, then a message from a friend.

\begin{itemize}
\raggedright
  \item \emph{Recall:} say \emph{huífù}, to reply — \textbf{R1}, one lesson back
  \item \emph{Recall:} say \emph{tōngzhī}, a notice — \textbf{R2}
\end{itemize}

\begin{quote}
\zh{大家}
\end{quote}

\emph{Read it:} it aloud, with nothing said first — \emph{dàjiā}
\end{quote}

\begin{tcolorbox}[breakable,skin=enhanced,colback=orange!6,colframe=orange!45!black,boxrule=0.5pt,arc=1mm,left=6pt,right=6pt,top=4pt,bottom=4pt,fonttitle=\bfseries,title={Read this}]
\begin{quote}
\zh{明天下午三点开会} \zh{大家都要来} \zh{请不要来晚了}
\end{quote}

\begin{quote}
\zh{我明天不能去了} \zh{因为我妈妈病了} \zh{后天下午你在家吗} \zh{我想两点来看你} \zh{给我打电话}
\end{quote}

\emph{Read it:} the notice and the message through once without stopping; do not translate as you go

\emph{Read it:} the message again, and find the question your friend asks you
\end{tcolorbox}

\begin{tcolorbox}[breakable,skin=enhanced,colback=orange!6,colframe=orange!45!black,boxrule=0.5pt,arc=1mm,left=6pt,right=6pt,top=4pt,bottom=4pt,fonttitle=\bfseries,title={Read this}]
\begin{quote}
\emph{míngtiān xiàwǔ sān diǎn kāihuì} \emph{dàjiā dōu yào lái} \emph{qǐng bú yào lái wǎn le}
\end{quote}

\begin{quote}
\emph{wǒ míngtiān bù néng qù le} \emph{yīnwèi wǒ māma bìng le} \emph{hòutiān xiàwǔ nǐ zài jiā ma} \emph{wǒ xiǎng liǎng diǎn lái kàn nǐ} \emph{gěi wǒ dǎ diànhuà}
\end{quote}

The notice speaks to everyone, \zh{大家}, and asks with \zh{请}. The message runs in order: what has changed, why, a question, an offer, a request. \zh{因为} ties the reason to the line before it.
\end{tcolorbox}

\subsection*{Your turn}
\begin{itemize}
\raggedright
  \item When is the meeting? (\textbf{\zh{明天下午三点}}: three tomorrow afternoon.)
  \item Why isn't your friend coming tomorrow? (\textbf{\zh{因为我妈妈病了}}: the mother is ill.)
  \item What does your friend ask you? (\textbf{\zh{后天下午你在家吗}}: whether you are home the afternoon after tomorrow.)
  \item What does your friend want you to do? (\textbf{\zh{给我打电话}}: phone back.)
\end{itemize}

\begin{itemize}
\raggedright
  \item \emph{Say it:} the message aloud, once, as your friend would say it
\end{itemize}

\begin{tcolorbox}[breakable,colback=teal!4,colframe=teal!35!black,title={Before you move on}]
Who is the notice for? (\textbf{\zh{大家}}: everyone.) What must nobody be? (\textbf{\zh{来晚了}}: late.) What time does your friend want to come? (\textbf{\zh{两点}}.)
\end{tcolorbox}
\section[(a reply)]{A reply, and the lines that lean on each other}

\label{lesson:ZH-W157-reply}



\begin{quote}
The last lesson ended on a message from a friend: they cannot come tomorrow, and they ask whether you are at home the afternoon after.

This lesson answers it.

\begin{itemize}
\raggedright
  \item \emph{Recall:} say \emph{huífù}, to reply
\end{itemize}
\end{quote}

\subsection*{What to know first}
Connected writing is not \textbf{longer} writing. Four short lines can hang together; six long ones can fall apart. What makes lines connected is that each one leans on the one before it.

Three things do the leaning, and you own all three. \textbf{A time word at the front} of a line tells your reader where in the day they are: \zh{后天下午}, \zh{两点以后}. \textbf{The subject said once}: say \zh{我} in the first line, and the lines after it can start on the time or the verb. \textbf{\zh{因为}} ties a reason to the line before it, so your reader knows why.

Your reader is the friend who sent the message. Your purpose is to settle the appointment: answer the question, then agree a time.

\subsection*{Writing — connected composition}
Write \textbf{four to six lines} back to your friend. No name at the top: write to them as you would in a message.

\begin{itemize}
\raggedright
  \item Line 1 answers the question: \zh{后天下午我在家}, or \zh{后天下午我不在家}.
  \item Line 2 says what you will be doing then, with \zh{会}, \zh{要} or \zh{正在}.
  \item Line 3 gives one reason, with \zh{因为}.
  \item Line 4 offers a time, with \zh{从} … \zh{到} …, or \zh{两点以后}.
  \item The last line asks after their mother. \zh{你妈妈} … needs only characters you write.
\end{itemize}

A word whose characters you have not written goes in pinyin, inside the line, with its tone marks: \emph{kěyǐ}, \emph{shūfu}, \emph{huífù}. Every character comes from the 163 you have learned, each in the order you learned it, each in the same size of square.

No Chinese commas or full stops: leave a space between clauses, as this book prints them.

Write the reply. Then read it once as your friend would, top to bottom.

\begin{tcolorbox}[breakable,colback=teal!4,colframe=teal!35!black,title={Before you move on}]
Read your third line on its own. Does it stand up without the line above it? (\textbf{It should not quite: \zh{因为} leans on the line before}.)

Count the lines that start with \zh{我}. If every one does, you wrote a list, not a reply.
\end{tcolorbox}
\section[(dialogue)]{The A2 core — once more}

\label{lesson:ZH-R157-review-final-157}



\begin{quote}
The last lesson of the A2 core: one conversation that uses every chapter of it.

\begin{itemize}
\raggedright
  \item \emph{Recall:} say \emph{xiāoxi}, a message — \textbf{R2}
  \item \emph{Recall:} say \zh{开门} and \emph{guānmén}, to open and to close — \textbf{R2}
  \item \emph{Recall:} say \emph{jiàqián} and \emph{miǎnfèi}, the price and free — \textbf{R2}
\end{itemize}
\end{quote}

\subsection*{Your turn}
\begin{quote}
\textbf{A:} \zh{明天会下雨吗}
\end{quote}

\begin{quote}
\emph{míngtiān huì xiàyǔ ma} — Will it rain tomorrow?
\end{quote}

\begin{quote}
\textbf{B:} \emph{tiānqì yùbào shuō kěnéng huì} — The forecast says it might.
\end{quote}

\begin{quote}
\textbf{B:} \zh{后天不会}
\end{quote}

\begin{quote}
\emph{hòutiān bú huì} — Not the day after.
\end{quote}

\begin{quote}
\textbf{A:} \zh{后天下午开会} \zh{你来吗}
\end{quote}

\begin{quote}
\emph{hòutiān xiàwǔ kāihuì, nǐ lái ma} — There's a meeting the afternoon after tomorrow. Are you coming?
\end{quote}

\begin{quote}
\textbf{B:} \emph{wǒ yídìng lái} — I'll be there, for sure.
\end{quote}

\begin{quote}
\textbf{A:} \zh{从两点到四点}
\end{quote}

\begin{quote}
\emph{cóng liǎng diǎn dào sì diǎn} — From two till four.
\end{quote}

\begin{quote}
\textbf{B:} \emph{nǐ kàndào tōngzhī le ma} — Did you see the notice?
\end{quote}

\begin{quote}
\textbf{A:} \emph{kàndào le, wǒ yě huífù le} — I saw it, and I've replied.
\end{quote}

\begin{quote}
\textbf{B:} \zh{下个月你要去中国吗}
\end{quote}

\begin{quote}
\emph{xià ge yuè nǐ yào qù Zhōngguó ma} — Are you going to China next month?
\end{quote}

\begin{quote}
\textbf{A:} \zh{要}
\end{quote}

\begin{quote}
\emph{yào} — I am.
\end{quote}

\begin{quote}
\textbf{A:} \zh{你去过中国吗}
\end{quote}

\begin{quote}
\emph{nǐ qù guo Zhōngguó ma} — Have you ever been?
\end{quote}

\begin{quote}
\textbf{B:} \zh{没去过}
\end{quote}

\begin{quote}
\emph{méi qù guo} — No, never.
\end{quote}

\begin{quote}
\textbf{B:} \zh{你汉语说得很好}
\end{quote}

\begin{quote}
\emph{nǐ Hànyǔ shuō de hěn hǎo} — Your Chinese is good.
\end{quote}

\begin{quote}
\textbf{A:} \zh{我没有他那么好}
\end{quote}

\begin{quote}
\emph{wǒ méiyǒu tā nàme hǎo} — Not as good as his.
\end{quote}

\begin{quote}
\textbf{B:} \zh{你要水还是要饭}
\end{quote}

\begin{quote}
\emph{nǐ yào shuǐ háishi yào fàn} — Water or food?
\end{quote}

\begin{quote}
\textbf{A:} \zh{饭} \zh{我还没吃}
\end{quote}

\begin{quote}
\emph{fàn, wǒ hái méi chī} — Food. I haven't eaten yet.
\end{quote}

\begin{quote}
\textbf{B:} \zh{我正在做饭}
\end{quote}

\begin{quote}
\emph{wǒ zhèngzài zuò fàn} — I'm cooking right now.
\end{quote}

\begin{quote}
\textbf{B:} \emph{bié zǒu} — Don't go.
\end{quote}

\begin{quote}
\textbf{A:} \emph{wǒ de huǒchē bā diǎn chūfā} — My train leaves at eight.
\end{quote}

\begin{quote}
\textbf{A:} \emph{wǒ kàndào yí ge guǎnggào} — I saw an advert:
\end{quote}

\begin{quote}
\textbf{A:} \zh{书十元一本}
\end{quote}

\begin{quote}
\emph{shū shí yuán yì běn} — books at ten yuan each.
\end{quote}

\begin{quote}
\textbf{B:} \zh{星期天不开门}
\end{quote}

\begin{quote}
\emph{xīngqītiān bù kāimén} — That shop is closed on Sundays.
\end{quote}

\begin{itemize}
\raggedright
  \item \emph{Read it:} the lines printed in characters aloud, then the rest
  \item \emph{Say it:} the whole dialogue, both parts
  \item \emph{Say it:} B's lines again, with your own answers
\end{itemize}

\begin{tcolorbox}[breakable,colback=teal!4,colframe=teal!35!black,title={Before you move on}]
Will? (\textbf{\zh{会}}.) Going to? (\textbf{\zh{要}}.) Right now? (\textbf{\zh{正在}}.) Ever? (\textbf{\zh{过}}.) Saw it? (\textbf{\zh{看到}}.) How well? (\textbf{\zh{得}}.) Than? (\textbf{\zh{比}}.) Or? (\textbf{\zh{还是}}.) Not yet? (\textbf{\zh{还没}}.) From … to? (\textbf{\zh{从} … \zh{到}}.) Yuan, as written? (\textbf{\zh{元}}.) A meeting? (\textbf{\zh{开会}}.) Reply? (\emph{huífù}.)
\end{tcolorbox}
